\section{Background}\label{sec:background}

\subsection{AlphaZero-style self-play}

Monte Carlo tree search~\citep{coulom2006efficient,kocsis2006bandit} guided by a policy
and value network is the basis of AlphaGo~\citep{silver2016mastering}, AlphaGo
Zero~\citep{silver2017mastering}, and AlphaZero~\citep{silver2018general}, together with
their open reproductions ELF OpenGo~\citep{tian2019elf}, Leela
Zero~\citep{pascutto2019leela} and SAI~\citep{morandin2019sai}.

Each playout descends from the root. At each node it picks the child $c$ that maximises
the PUCT score~\citep{rosin2011multi,silver2017mastering}:
\begin{equation}
  \mathrm{PUCT}(c) = V(c) + c_{\mathrm{PUCT}}\,P(c)\,\frac{\sqrt{\sum_{c'} N(c')}}{1 + N(c)},
  \label{eq:puct}
\end{equation}
where $V$ is the mean value of the subtree, $P$ the network's prior and $N$ the visit
count. The root visit distribution $\pi$ becomes the policy target, the game outcome $z$
the value target, and the network is trained to minimise
\begin{equation}
  L = (z - \hat v)^2 - \textstyle\sum_m \pi(m)\log\hat\pi(m) + c\,\lVert\theta\rVert^2 .
  \label{eq:az-loss}
\end{equation}
AlphaGo Zero also exploited the eight rotations and reflections of the board to augment
its training data~\citep{silver2017mastering,silver2018general}, and it gated each new
network by requiring a 55\% win rate against the current one over 400
games~\citep{silver2017mastering}. AlphaZero dropped gating and trained
continuously~\citep{silver2018general}.

\subsection{The \katago{} recipe}\label{sec:bg-katago}

\citet{wu2019accelerating} keeps the AlphaZero skeleton and changes the parts listed
below. This report refers back to them by name.
\begin{description}
  \item[Playout cap randomization.] Only a fraction $p=0.25$ of turns gets a full search
    (a cap of $N$ nodes); the rest use a fast search with cap $n<N$. Early in the run
    $(N,n)=(600,100)$, rising to $(1000,200)$. \emph{Only full-search turns are recorded
    for training}. Fast turns exist to generate more games, and therefore more value
    targets, per unit of compute. Informal evidence from Leela Zero suggests that policy
    learning needs far more playouts per move than value learning~\citep{forsten2019optimal}.
  \item[Forced playouts and policy target pruning.] Children of the root get a minimum
    number of playouts, $n_{\mathrm{forced}}(c)=(k P(c)\sum N)^{1/2}$ with $k=2$. The
    forced extras are then subtracted from the policy target, so exploration no longer
    pollutes what the policy learns.
  \item[Global pooling.] Pooling layers compute each channel's mean, its mean scaled by board
    width, and its maximum. A fully connected layer maps these to a per-channel bias for
    another set of channels. The structure follows the first convolution of some residual
    blocks and appears in the policy and value heads. It is related to
    Squeeze-and-Excitation~\citep{hu2018squeeze}.
  \item[Auxiliary policy target.] The policy head also predicts the opponent's reply, with
    weight $w_{\mathrm{opp}}=0.15$~\citep{tian2016better}.
  \item[Ownership and score targets.] The network predicts the final owner of each point
    (cross-entropy with $w_o = 1.5/b^2$) and the final score, via a ``pdf'' and a ``cdf''
    loss each with weight $0.02$. This follows the multi-labelled value networks of
    \citet{wu2018multilabelled}.
  \item[Game-specific input features.] Liberties, ladders, pass-alive regions and komi
    parity.
  \item[Training details.] SGD with momentum 0.9, batch 256, $c_{L2}=3\times10^{-5}$,
    value-loss scale $c_g=1.5$. The per-sample learning rate is fixed, lowered for the
    first 5M samples and near the end. Samples are drawn uniformly from a moving window
    whose size grows sublinearly with the total amount of data:
    \begin{equation}
      N_{\mathrm{window}} = c\left(1 + \beta\,\frac{(N_{\mathrm{total}}/c)^{\alpha} - 1}{\alpha}\right),
      \quad c=250{,}000,\ \alpha=0.75,\ \beta=0.4
      \label{eq:window}
    \end{equation}
    \citep[App.~C]{wu2019accelerating}. Candidates are an exponential moving average of
    weight snapshots, in the spirit of stochastic weight
    averaging~\citep{izmailov2018averaging}. A candidate must win at least 100 of 200 games
    to replace the current network~\citep[App.~E]{wu2019accelerating}.
\end{description}
\citeauthor{wu2019accelerating}'s ablations (Table~2 of that paper) put the training-time
factor of each technique between $1.25\times$ and $1.65\times$ at about two days of
compute. The factors multiply to roughly $9\times$. This report uses those factors only as
motivation. We do not reproduce them.

\subsection{\gofer{} before this work}\label{sec:bg-gofer}

\gofer{} plays under Chinese area scoring with Tromp--Taylor-style
evaluation~\citep{tromp2014go}. Its network interface is fixed by a feature schema:
\begin{itemize}
  \item \textbf{Spatial input} ($8\times S\times S$): own stones, opponent stones, empty
        points, the simple-ko point, a to-move plane, and the last three moves.
  \item \textbf{Global input} (4 values): komi$/10$, the normalised move number, and one-hot
        side to move.
  \item \textbf{Outputs}: $S^2+1$ policy logits (pass last) and a scalar $\tanh$ value from the
        side to move's perspective.
\end{itemize}

Search uses \cref{eq:puct} with $c_{\mathrm{PUCT}}=1.1$. Unvisited children get a fixed
value of $-0.2$, whereas \katago{}'s first-play urgency is the parent value minus
$0.2\sqrt{P_{\mathrm{explored}}}$. In self-play, root Dirichlet noise ($\alpha=0.15$, blend $0.25$) is applied on
every move, fast and full, whereas \katago{} disables it on fast searches. Full searches
also force a minimum number of visits on each root child ($2 + \lfloor\sqrt{P(c)(N+1)}\rfloor$
for playout cap $N$, versus \katago{}'s $\sqrt{2P(c)\sum N}$), and they prune the policy
target by dropping children with fewer than two visits (\cref{sec:selfplay-search}). Inference runs through ONNX
Runtime~\citep{onnxruntime2021}, either in-process or through a Python sidecar.

The v3 loop, documented in the repository's ADR~0003, ran the following cycle:
\begin{enumerate}
  \item Self-play: 200 games per cycle, 200 playouts per move with playout cap randomization
        ($p=0.20$, fast cap 50), keeping only full-search rows.
  \item Replay: append the games to a JSONL replay file capped first-in-first-out at 50{,}000
        rows.
  \item Training: fine-tune the champion's 6-block, 64-channel network (``6$\times$64'') for
        15 epochs with SGD.
  \item Gating: export the candidate to ONNX and play up to 200 arena games against the
        champion. Promote if it scores at least 0.55 with the lower Wilson
        bound~\citep{wilson1927probable} above 0.5. The arena stops early: after 20 games
        once 0.55 is out of reach, and after 100 games once the promotion condition holds.
\end{enumerate}

ADR~0005 compared net sizes offline and found a thin validation-loss edge for a 4$\times$48
network. It kept 6$\times$64 anyway, because changing architecture breaks warm-starting and
was estimated to cost 4--5 days and about 20 cycles to recover the champion's strength on
the project's CPU instance. \Cref{sec:learner-distill} targets exactly that cost.
